% !TEX program = pdflatex
% ==========================================================
% Universidad Santo Tomás
% Fundamentos de Matemática Básica
% Clase 3: Potencias
% Luis Tulio González Alcaino
% ==========================================================

\newif\ifhandout
\handoutfalse
%\handouttrue

\ifhandout
\documentclass[aspectratio=169,10pt,handout]{beamer}
\else
\documentclass[aspectratio=169,10pt]{beamer}
\fi

% ----------------------------------------------------------
% PAQUETES
% ----------------------------------------------------------
\usepackage[spanish,es-noshorthands]{babel}
\usepackage[T1]{fontenc}
\usepackage[utf8]{inputenc}

\usepackage{lmodern}
\usepackage{amsmath,amsfonts,amssymb}
\usepackage{ragged2e}
\usepackage{enumitem}
\usepackage{xcolor}
\usepackage{tikz}
\usetikzlibrary{positioning}

% ----------------------------------------------------------
% COLORES INSTITUCIONALES
% ----------------------------------------------------------
\definecolor{USTBlue}{RGB}{0,70,127}
\definecolor{USTBlueDark}{RGB}{0,45,85}
\definecolor{USTTeal}{RGB}{0,153,117}
\definecolor{USTGray}{RGB}{120,120,120}
\definecolor{USTSoft}{RGB}{248,249,250}
\definecolor{USTText}{RGB}{35,40,45}
\definecolor{WarmGray}{RGB}{120,120,120}

% ----------------------------------------------------------
% TEMA BASE
% ----------------------------------------------------------
\usetheme{Madrid}
\usecolortheme[named=USTBlue]{structure}
\setbeamertemplate{navigation symbols}{}

% ----------------------------------------------------------
% AJUSTES VISUALES
% ----------------------------------------------------------
\setlength{\parskip}{0.2em}
\setlist[itemize]{itemsep=0.35em}
\setlist[enumerate]{itemsep=0.35em}

% ----------------------------------------------------------
% COLORES DE PIE
% ----------------------------------------------------------
\setbeamercolor{author in head/foot}{bg=USTSoft,fg=USTGray}
\setbeamercolor{date in head/foot}{bg=USTSoft,fg=USTGray}
\setbeamercolor{title in head/foot}{bg=USTSoft,fg=USTGray}

% ----------------------------------------------------------
% PIE DE PÁGINA
% ----------------------------------------------------------
\setbeamertemplate{footline}{%
	\leavevmode%
	\hbox{%
		\begin{beamercolorbox}[wd=.85\paperwidth,ht=2.5ex,dp=1ex,leftskip=0.4cm]{author in head/foot}%
			\scriptsize
			\textbf{Fundamentos de Matemática Básica}
			\quad|\quad
			Agronomía
			\quad|\quad
			Universidad Santo Tomás Talca
		\end{beamercolorbox}%
		\begin{beamercolorbox}[wd=.15\paperwidth,ht=2.5ex,dp=1ex,rightskip=0.3cm plus1fil]{date in head/foot}%
			\scriptsize \insertframenumber/\inserttotalframenumber
		\end{beamercolorbox}%
	}%
}

% ----------------------------------------------------------
% BLOQUES PEDAGÓGICOS
% ----------------------------------------------------------
\setbeamercolor{block title}{bg=USTBlue!15,fg=USTBlueDark}
\setbeamercolor{block body}{bg=USTSoft,fg=black}

\setbeamercolor{block title example}{bg=USTTeal!20,fg=USTBlueDark}
\setbeamercolor{block body example}{bg=USTSoft,fg=black}

\setbeamercolor{block title alerted}{bg=red!15,fg=black}
\setbeamercolor{block body alerted}{bg=red!3,fg=black}

% ----------------------------------------------------------
% COMANDOS Y ENTORNOS AUXILIARES
% ----------------------------------------------------------
\newcommand{\destacar}[1]{\textcolor{USTBlueDark}{\textbf{#1}}}
\newcommand{\alerta}[1]{\textcolor{red!70!black}{\textbf{#1}}}
\newcommand{\pp}{\pause}

\newenvironment{metaBox}
{\begin{block}{Meta de aprendizaje}}
	{\end{block}}

\newenvironment{ideaBox}
{\begin{exampleblock}{Idea clave}}
	{\end{exampleblock}}

% ----------------------------------------------------------
% DATOS DEL CURSO
% ----------------------------------------------------------
\title{Fundamentos de Matemática Básica}
\subtitle{Unidad I: Operaciones algebraicas en los reales\\Clase 3: Potencias}
\author{Prof. Luis González Alcaino}
\date{Universidad Santo Tomás -- Agronomía 2026}

% ==========================================================
% DOCUMENTO
% ==========================================================
\begin{document}
	
	% ----------------------------------------------------------
	\begin{frame}
		\titlepage
	\end{frame}
	
	% ----------------------------------------------------------
	\begin{frame}{Contenidos de la clase}
		\begin{block}{Aprendizajes de la sesión}
			\begin{enumerate}
				\item Comprender el concepto de \destacar{potencia}.
				\item Interpretar la \destacar{base} y el \destacar{exponente}.
				\item Aplicar las \destacar{propiedades de las potencias}.
				\item Reconocer \destacar{patrones numéricos}.
				\item Resolver \destacar{problemas contextualizados en Agronomía}.
			\end{enumerate}
		\end{block}
	\end{frame}
	
	% ----------------------------------------------------------
	
	% ----------------------------------------------------------
	\begin{frame}{¿Por qué estudiar potencias en Agronomía?}
		\begin{columns}[T,totalwidth=\textwidth]
			\begin{column}{0.52\textwidth}
				\begin{block}{Situaciones reales}
					\begin{itemize}
						\item Número de semillas por hectárea.
						\item Cantidad de bacterias en una muestra de suelo.
						\item Concentración de micronutrientes.
						\item Conversión de unidades y escalas.
					\end{itemize}
				\end{block}
			\end{column}
			
			\begin{column}{0.44\textwidth}
				\begin{exampleblock}{Ejemplo rápido}
					Una dosis puede expresarse como:
					\[
					0{,}00045\ \text{kg}=4{,}5\times10^{-4}\ \text{kg}
					\]
					La segunda forma es más clara para calcular, comparar y comunicar.
				\end{exampleblock}
			\end{column}
		\end{columns}
		
		\vspace{0.2cm}
		\centering
		{\small \textcolor{WarmGray}{La matemática permite simplificar, modelar y validar información técnica.}}
	\end{frame}
	
	% ----------------------------------------------------------
	\begin{frame}{Concepto de potencia}
		\begin{block}{Definición}
			Una \destacar{potencia} representa una multiplicación repetida de un mismo factor:
			\[
			a^n=\underbrace{a\cdot a\cdot a\cdots a}_{n \text{ veces}}
			\]
			donde:
			\begin{itemize}
				\item \(a\) es la \destacar{base},
				\item \(n\) es el \destacar{exponente}.
			\end{itemize}
		\end{block}
		
		\pp
		
		\begin{exampleblock}{Ejemplos}
			\[
			2^4 = 2\cdot2\cdot2\cdot2 = 16
			\qquad
			10^3 = 1000
			\]
			\[
			5^2 = 25
			\qquad
			10^{-2}=\frac{1}{10^2}=\frac{1}{100}=0{,}01
			\]
		\end{exampleblock}
	\end{frame}
	
	% ----------------------------------------------------------
	\begin{frame}{Propiedades de las potencias I}
		\begin{block}{1. Producto de potencias de igual base}
			\[
			a^m \cdot a^n = a^{m+n}
			\]
		\end{block}
		
		\begin{exampleblock}{Ejemplo}
			\[
			2^3\cdot2^4 = 2^{3+4}=2^7=128
			\]
		\end{exampleblock}
		
		\pause
		
		\begin{block}{2. Cociente de potencias de igual base}
			\[
			\frac{a^m}{a^n}=a^{m-n}, \qquad a\neq 0
			\]
		\end{block}
		
		\begin{exampleblock}{Ejemplo}
			\[
			\frac{5^6}{5^2}=5^{6-2}=5^4=625
			\]
		\end{exampleblock}
	\end{frame}
	
	% ----------------------------------------------------------
	\begin{frame}{Propiedades de las potencias II}
		\begin{block}{3. Potencia de una potencia}
			\[
			(a^m)^n = a^{m\cdot n}
			\]
		\end{block}
		
		\begin{exampleblock}{Ejemplo}
			\[
			(3^2)^4 = 3^{2\cdot4}=3^8=6561
			\]
		\end{exampleblock}
		
		\pause
		
		\begin{block}{4. Potencia de un producto}
			\[
			(ab)^n = a^n b^n
			\]
		\end{block}
		
		\begin{exampleblock}{Ejemplo}
			\[
			(2\cdot5)^3 = 2^3\cdot5^3 = 8\cdot125 = 1000
			\]
		\end{exampleblock}
	\end{frame}
	
	% ----------------------------------------------------------
	\begin{frame}{Propiedades de las potencias III}
		\begin{block}{5. Potencia de un cociente}
			\[
			\left(\frac{a}{b}\right)^n = \frac{a^n}{b^n}, \qquad b\neq 0
			\]
		\end{block}
		
		\begin{exampleblock}{Ejemplo}
			\[
			\left(\frac{2}{3}\right)^3 = \frac{2^3}{3^3}=\frac{8}{27}
			\]
		\end{exampleblock}
		
		\pause
		
		\begin{block}{6. Exponente cero}
			\[
			a^0 = 1, \qquad a\neq 0
			\]
		\end{block}
		
		\begin{exampleblock}{Ejemplo}
			\[
			7^0 = 1
			\]
		\end{exampleblock}
	\end{frame}
	
	% ----------------------------------------------------------
	\begin{frame}{Propiedades de las potencias IV}
		\begin{block}{7. Exponente uno}
			\[
			a^1 = a
			\]
		\end{block}
		
		\begin{exampleblock}{Ejemplo}
			\[
			9^1 = 9
			\]
		\end{exampleblock}
		
		\pause
		
		\begin{block}{8. Exponente negativo}
			\[
			a^{-n}=\frac{1}{a^n}, \qquad a\neq 0
			\]
		\end{block}
		
		\begin{exampleblock}{Ejemplo}
			\[
			2^{-3}=\frac{1}{2^3}=\frac{1}{8}
			\]
		\end{exampleblock}
	\end{frame}
	
	% ----------------------------------------------------------
	\begin{frame}{Propiedades de las potencias V}
		\begin{block}{9. Base igual a uno}
			\[
			1^n = 1
			\]
		\end{block}
		
		\begin{exampleblock}{Ejemplo}
			\[
			1^8 = 1
			\]
		\end{exampleblock}
		
		\pause
		
		\begin{block}{10. Base igual a cero}
			\[
			0^n = 0, \qquad n>0
			\]
		\end{block}
		
		\begin{exampleblock}{Ejemplo}
			\[
			0^5 = 0
			\]
		\end{exampleblock}
	\end{frame}
	
	% ----------------------------------------------------------
	\begin{frame}{Importante: errores frecuentes}
		\begin{alertblock}{Cuidado}
			No se pueden sumar ni restar exponentes en una suma de potencias:
			\[
			a^m + a^n \neq a^{m+n}
			\]
		\end{alertblock}
		
		\pause
		
		\begin{exampleblock}{Ejemplo}
			\[
			2^2 + 2^3 = 4 + 8 = 12
			\]
			pero
			\[
			2^{2+3}=2^5=32
			\]
			Por lo tanto,
			\[
			2^2 + 2^3 \neq 2^5
			\]
		\end{exampleblock}
	\end{frame}
	
	% ----------------------------------------------------------
	\begin{frame}{Patrones}
		Observe la secuencia:
		\[
		2^1 = 2
		\]
		\[
		2^2 = 4
		\]
		\[
		2^3 = 8
		\]
		\[
		2^4 = 16
		\]
		
		\begin{block}{Preguntas}
			\begin{itemize}
				\item ¿Qué ocurre cuando el exponente aumenta en 1?
				\item ¿Cómo crecerá la secuencia?
			\end{itemize}
		\end{block}
	\end{frame}
	
	% ----------------------------------------------------------
	\begin{frame}{Patrones y relaciones en las potencias de 10}
		\begin{columns}[T,totalwidth=\textwidth]
			\begin{column}{0.48\textwidth}
				\begin{block}{Observación}
					\[
					10^1 = 10
					\]
					\[
					10^2 = 100
					\]
					\[
					10^3 = 1000
					\]
					\[
					10^4 = 10000
					\]
					Cada vez que el exponente aumenta en 1, el número se multiplica por 10.
				\end{block}
			\end{column}
			
			\begin{column}{0.48\textwidth}
				\begin{exampleblock}{Hacia exponentes negativos}
					\[
					10^0 = 1
					\]
					\[
					10^{-1} = 0{,}1
					\]
					\[
					10^{-2} = 0{,}01
					\]
					\[
					10^{-3} = 0{,}001
					\]
					Cada vez que el exponente disminuye en 1, el número se divide por 10.
				\end{exampleblock}
			\end{column}
		\end{columns}
		
		\vspace{0.2cm}
		\begin{ideaBox}
			\textbf{Patrón clave:} los exponentes permiten describir regularidades de crecimiento o disminución en escala decimal.
		\end{ideaBox}
	\end{frame}
	
	% ----------------------------------------------------------
	\begin{frame}{Ejercicios}
		Resuelva:
		\begin{enumerate}
			\item \(2^4 \cdot 2^3\)
			\item \((3^2)^3\)
			\item \(5^3 \cdot 5^2\)
			\item \(\dfrac{10^6}{10^2}\)
			\item En un cultivo se plantan \(3^4\) semillas por parcela. Si hay \(3^2\) parcelas, ¿cuántas semillas hay en total?
		\end{enumerate}
	\end{frame}
	
	% ----------------------------------------------------------
	\begin{frame}{Ejercicio contextualizado}
		\begin{block}{Situación agrícola}
			Una parcela experimental utiliza:
			\[
			2^5
			\]
			semillas por metro cuadrado.
			
			Si el cultivo cubre:
			\[
			2^3
			\]
			metros cuadrados, ¿cuántas semillas se utilizan en total?
		\end{block}
		
		\pause
		
		\begin{exampleblock}{Resolución}
			\[
			2^5 \cdot 2^3 = 2^{5+3} = 2^8
			\]
			\[
			2^8 = 256
			\]
			Se utilizan aproximadamente \textbf{256 semillas}.
		\end{exampleblock}
	\end{frame}
	
	% ----------------------------------------------------------
	\begin{frame}{Crecimiento de bacterias en el suelo}
		\begin{block}{Contexto}
			Una colonia bacteriana se duplica cada hora.
			
			La cantidad inicial es:
			\[
			2^3
			\]
			bacterias.
			
			Después de 4 horas:
			\[
			2^3 \cdot 2^4
			\]
		\end{block}
		
		\pause
		
		\begin{exampleblock}{Cálculo}
			\[
			2^3 \cdot 2^4 = 2^{7}
			\]
			\[
			2^7 = 128
			\]
			Después de 4 horas hay \textbf{128 bacterias}.
		\end{exampleblock}
	\end{frame}
	
	% ----------------------------------------------------------
	\begin{frame}{Fertilización}
		\begin{block}{Situación}
			Un fertilizante se aplica en dosis proporcionales a potencias de 10.
			
			En un ensayo se utiliza:
			\[
			5\times10^2
			\]
			gramos por hectárea.
			
			Si el cultivo ocupa:
			\[
			10^3
			\]
			hectáreas experimentales, calcule la cantidad total.
		\end{block}
		
		\pause
		
		\begin{exampleblock}{Resolución}
			\[
			(5\times10^2)(10^3)=5\times10^{5}
			\]
			Se requieren aproximadamente:
			\[
			500000\text{ gramos}
			\]
		\end{exampleblock}
	\end{frame}
	
	% ----------------------------------------------------------
	\begin{frame}{Ejercicios contextualizados}
		\begin{enumerate}
			\item En una parcela experimental se plantan
			\[
			2^6
			\]
			semillas de trigo por sector.
			
			Si el terreno tiene
			\[
			2^3
			\]
			sectores iguales, ¿cuántas semillas se plantan en total?
			
			\pause
			
			\item Una bacteria del suelo se duplica cada hora. Si inicialmente hay
			\[
			2^2
			\]
			bacterias, ¿cuántas habrá después de 6 horas?
			
			\pause
			
			\item Un cultivo utiliza una densidad de
			\[
			5^3
			\]
			plantas por unidad experimental.
			
			Si se instalan
			\[
			5^2
			\]
			unidades experimentales, ¿cuántas plantas se utilizan en total?
			
			\pause
			
			\item En un laboratorio agrícola se almacenan muestras de suelo. Cada caja contiene
			\[
			3^4
			\]
			muestras.
			
			Si se tienen
			\[
			3^2
			\]
			cajas, ¿cuántas muestras hay en total?
		\end{enumerate}
	\end{frame}
	
	% ----------------------------------------------------------
	\begin{frame}{Síntesis}
		\begin{block}{Ideas clave}
			\begin{itemize}
				\item Las potencias representan multiplicaciones repetidas.
				\item Las propiedades de las potencias permiten simplificar cálculos.
				\item Las potencias aparecen frecuentemente en fenómenos de crecimiento.
				\item En Agronomía permiten modelar poblaciones, concentraciones y escalas.
			\end{itemize}
		\end{block}
	\end{frame}
	
		\begin{frame}{Trabajo personal del estudiante}
		
		\centering
		
		\vspace{0.3cm}
		
		Libro Álgebra, trigonometría y geometría análitica
		Dennis G. Zill - Jaqueline.M Dewar
		
		\medskip
		
		Ejercicio 2.3 Páginas 68-70	
		
	\end{frame}	
	
	
\end{document}